\documentclass[11pt,a4paper]{article}

\usepackage[utf8]{inputenc}
\usepackage{geometry}
\geometry{
    a4paper,
    top=2.5cm,
    bottom=2.5cm,
    left=2.5cm,
    right=2.5cm
}
\usepackage{hyperref}
\usepackage{amsmath}
\usepackage{enumitem}
\usepackage{titlesec}
\usepackage{xcolor}
\usepackage{calc}
\usepackage[most]{tcolorbox}

% Style configuration
\definecolor{primary}{RGB}{31, 78, 121}

% Timeline Box Style
\newtcolorbox{timelinebox}[2][]{
    enhanced,
    boxrule=0pt,
    leftrule=3.5pt,
    colframe=primary,
    colback=primary!4!white,
    sharp corners,
    fonttitle=\bfseries\color{primary}\large,
    title={#2},
    attach title to upper,
    after title={\par\smallskip},
    left=12pt,
    right=12pt,
    top=8pt,
    bottom=8pt,
    before=\noindent,
    after=\vspace{10pt},
    #1
}
\hypersetup{
    colorlinks=true,
    linkcolor=primary,
    filecolor=primary,      
    urlcolor=primary,
    citecolor=primary,
}

\titleformat{\section}
  {\normalfont\Large\bfseries\color{primary}}{\thesection}{1em}{}
\titleformat{\subsection}
  {\normalfont\large\bfseries\color{primary}}{\thesubsection}{1em}{}

\title{\textbf{\color{primary}Proposal: Dynamic Hybrid Retrieval-Augmented Generation for Vietnamese}}
\author{\textbf{Project Proposal}}
\date{\today}

\begin{document}

\maketitle

\section{Project Objective}
The objective of this project is to develop \textbf{Dynamic Hybrid RAG}, a lightweight framework designed specifically for Vietnamese Retrieval-Augmented Generation. Instead of relying on a one-size-fits-all combination of retrieval signals (dense vector search, sparse lexical matching like BM25, and learned-sparse representations), this project proposes an adaptive Multi-Layer Perceptron (MLP). This MLP will dynamically predict per-query three-way fusion weights based on Vietnamese-specific linguistic features.

\section{Key Findings \& Motivations}
From initial research and empirical analysis, several key insights motivate this work:
\begin{itemize}[leftmargin=*]
    \item \textbf{Vietnamese Specific Challenges:} Retrieval effectiveness in Vietnamese is highly sensitive to three main factors: word segmentation dependency (words spanning multiple syllables), diacritical mark presence (users often type without tone marks), and frequent code-switching with English (especially in technical domains).
    \item \textbf{Signal Complementarity:} Classical BM25 heavily relies on proper word segmentation and fails drastically when diacritics are missing. Dense embeddings are robust to diacritic noise but can blur specific lexical details. Learned-sparse representations (like BGE-M3) offer a middle ground, preserving token-level importance and handling code-switched terms better than BM25.
    \item \textbf{Superiority of Dynamic Fusion:} An adaptive weighting strategy significantly outperforms any single-signal retriever (dense-only, BM25-only, or sparse-only) as well as fixed-weight hybrid baselines.
    \item \textbf{Domain Generalization:} A dynamic MLP trained on a specific domain (e.g., Wikipedia via UIT-ViQuAD 2.0) can successfully generalize to unseen domains (e.g., legal/administrative texts in DANGDOCAO) in a zero-shot setting, proving it learns domain-invariant linguistic signals rather than dataset-specific lexical statistics.
\end{itemize}

\section{Implementation Methodology}
The implementation will be carried out through the following technical components:
\begin{itemize}[leftmargin=*]
    \item \textbf{Retrieval Components:}
    \begin{itemize}
        \item \emph{Dense Retrieval:} Using FPT Vietnamese Embedding (1024-dim) indexed with FAISS for semantic search.
        \item \emph{BM25 Retrieval:} Using \texttt{underthesea} for Vietnamese word segmentation and \texttt{BM25Okapi} for lexical matching.
        \item \emph{Learned Sparse Retrieval:} Using BGE-M3 lexical weights via an inverted index to capture learned term importance.
    \end{itemize}
    \item \textbf{Feature Extraction:} A dedicated module to extract 8 lightweight linguistic features from the query: diacritic ratio, compound word ratio, English token ratio, tech-term ratio, clause count, question-word presence, normalized query length, and out-of-vocabulary ratio against the BM25 corpus.
    \item \textbf{Adaptive MLP Fusion:} A 3-layer feed-forward neural network (Keras/TensorFlow, approx. 2,947 parameters; Dense+LayerNorm+GELU+Dropout blocks) taking the 8 features as input and outputting a softmax distribution over the 3 retrieval signals (dense, BM25, sparse) to guarantee weights sum to 1.
    \item \textbf{Soft-Label Training Strategy:} To train the MLP without ground-truth weights, we will compute NDCG@10 across a 3D simplex grid of possible weights (231 points at step 0.05). We will then apply a temperature-scaled softmax to create smooth, expected-weight targets, minimizing MSE loss against these expected targets.
\end{itemize}

\section{Contributions}
While hybrid retrieval is well-studied for English, our approach introduces several novelties:
\begin{itemize}[leftmargin=*]
    \item \textbf{Dynamic vs. Fixed Weights:} Unlike standard approaches (e.g., Reciprocal Rank Fusion or fixed linear interpolation) that use the same static weights for all queries, our system computes custom fusion weights \emph{per query}.
    \item \textbf{Vietnamese-Aware Design:} The system explicitly models the nuances of the Vietnamese language through a custom 8-feature extractor, rather than applying language-agnostic statistical methods.
    \item \textbf{Soft-Label Simplex Supervision:} Instead of using hard argmax labels (which cause the model to collapse onto a single retriever and perform poorly), we introduce a temperature-scaled soft-label strategy over a 3D simplex grid, avoiding tie-breaking ambiguity and ensuring smooth gradient updates.
\end{itemize}

\section{Timeline}

\begin{timelinebox}{Week 1: Project Setup \& Data Preparation}
\begin{itemize}[leftmargin=*,nosep]
    \item Literature review and baseline setup.
    \item Acquire and preprocess datasets: UIT-ViQuAD 2.0 (Wikipedia) and DANGDOCAO (Legal/Administrative).
\end{itemize}
\end{timelinebox}

\begin{timelinebox}{Week 2: Implementation of Base Retrievers}
\begin{itemize}[leftmargin=*,nosep]
    \item Set up Dense retrieval (FPT embeddings + FAISS).
    \item Set up BM25 retrieval (\texttt{underthesea} segmentation + \texttt{BM25Okapi}).
    \item Set up BGE-M3 Learned Sparse retrieval (inverted index).
\end{itemize}
\end{timelinebox}

\begin{timelinebox}{Week 3: Feature Extractor Development}
\begin{itemize}[leftmargin=*,nosep]
    \item Implement the Vietnamese-aware 8-feature extractor.
    \item Run feature extraction over the training and dev query sets.
\end{itemize}
\end{timelinebox}

\begin{timelinebox}{Week 4: Fusion Model \& Training Pipeline}
\begin{itemize}[leftmargin=*,nosep]
    \item Build the 3-layer MLP fusion module.
    \item Implement the 3D simplex grid search and the soft-label supervision target generator.
\end{itemize}
\end{timelinebox}

\begin{timelinebox}{Week 5: Training \& In-Domain Evaluation}
\begin{itemize}[leftmargin=*,nosep]
    \item Train the MLP on augmented UIT-ViQuAD 2.0 data (including diacritic-stripped queries).
    \item Evaluate NDCG, MRR, MAP, and Recall against fixed baselines on the in-domain test set.
\end{itemize}
\end{timelinebox}

\begin{timelinebox}{Week 6: Cross-Domain \& Robustness Testing}
\begin{itemize}[leftmargin=*,nosep]
    \item Conduct zero-shot cross-domain evaluation on the DANGDOCAO corpus.
    \item Perform robustness analysis using diacritic-removed (noisy) query sets.
\end{itemize}
\end{timelinebox}

\begin{timelinebox}{Week 7: End-to-End RAG Evaluation}
\begin{itemize}[leftmargin=*,nosep]
    \item Integrate the retrieval pipeline with a generator (e.g., Qwen3-32B).
    \item Evaluate end-to-end QA quality using RAGAS metrics (Context Precision, Context Recall, Faithfulness, Answer Relevancy).
\end{itemize}
\end{timelinebox}

\begin{timelinebox}{Week 8: Finalization \& Documentation}
\begin{itemize}[leftmargin=*,nosep]
    \item Conduct ablation studies (e.g., varying the soft-label temperature).
    \item Finalize code refactoring, synthesize results, and complete the final research report/paper.
\end{itemize}
\end{timelinebox}

\end{document}
